\textbf{Lösung zu Klausur 29.07.2026, Aufgabe 5}

\teil{a} Kontrollmatrix (Stellen $1, \dots, 10$):
\[
H = \begin{pmatrix}
1 & 1 & 1 & 1 & 1 & 1 & 1 & 1 & 1 & 1\\
1 & 2 & 3 & 4 & 5 & 6 & 7 & 8 & 9 & 10
\end{pmatrix}
\]
Ein Wort $c$ ist Codewort genau dann, wenn $H c^T = 0$, also
$\sum_{i=1}^{10} c_i \equiv 0$ und $\sum_{i=1}^{10} i \cdot c_i \equiv 0 \pmod{11}$.

\teil{b} Kontrollziffern für $a = 1\,1\,5\,0\,0\,0\,0\,5$ ($c_i = a_i$ für $i = 1, \dots, 8$):
\begin{align*}
c_9 &\equiv \sum_{i=1}^{8} (1+i)\, c_i = 2 \cdot 1 + 3 \cdot 1 + 4 \cdot 5 + 9 \cdot 5 = 70 \equiv 4 \pmod{11},\\
c_{10} &\equiv \sum_{i=1}^{8} (9-i)\, c_i = 8 \cdot 1 + 7 \cdot 1 + 6 \cdot 5 + 1 \cdot 5 = 50 \equiv 6 \pmod{11}.
\end{align*}
\begin{ergebnis} $c_9 = 4$, $c_{10} = 6$, $c = 1\,1\,5\,0\,0\,0\,0\,5\,4\,6$ \end{ergebnis}

\teil{c} Syndrome $S(r) = H r^T = (s_1, s_2)^T$:
\begin{align*}
r_1 = 1\,1\,5\,0\,0\,0\,0\,5\,1\,1:\quad
s_1 &= 1 + 1 + 5 + 5 + 1 + 1 = 14 \equiv 3,\\
s_2 &= 1 + 2 + 15 + 40 + 9 + 10 = 77 \equiv 0 \quad\Rightarrow\quad S(r_1) = (3, 0)^T,\\[4pt]
r_2 = 1\,1\,5\,0\,0\,0\,0\,7\,3\,3:\quad
s_1 &= 1 + 1 + 5 + 7 + 3 + 3 = 20 \equiv 9,\\
s_2 &= 1 + 2 + 15 + 56 + 27 + 30 = 131 \equiv 10 \quad\Rightarrow\quad S(r_2) = (9, 10)^T .
\end{align*}

\teil{d} Bei genau einem Fehler der Größe $e_x$ an der Stelle $x$ gilt $S = e_x \cdot (1, x)^T$,
also $e_x = s_1$ und $x = s_2 / s_1$.

Für $r_1$: $x = 0 / 3 = 0$. Die Stelle $0$ gibt es nicht (Stellen $1, \dots, 10$).
Das Syndrom passt also zu keinem Einzelfehler: $r_1$ enthält mindestens zwei Fehler.
Der Code hat $d = 3$ und korrigiert nur $1$ Fehler.
\begin{ergebnis} $r_1$ ist nicht dekodierbar (mindestens 2 Fehler). \end{ergebnis}

\teil{e} Für $r_2$: $e_x = s_1 = 9$ und $x = s_2 / s_1 = 10 \cdot 9^{-1} = 10 \cdot 5 = 50 \equiv 6 \pmod{11}$
(denn $9 \cdot 5 = 45 \equiv 1$). Fehler an Stelle $6$ mit Größe $9$:
\[
c_6 = r_6 - 9 = 0 - 9 \equiv 2 \pmod{11}.
\]
\begin{ergebnis} $c = 1\,1\,5\,0\,0\,2\,0\,7\,3\,3$ \end{ergebnis}

\teil{f} Probe $H c^T = 0$:
\begin{align*}
\text{Zeile 1:}\quad & 1 + 1 + 5 + 2 + 7 + 3 + 3 = 22 \equiv 0 \pmod{11},\\
\text{Zeile 2:}\quad & 1 + 2 + 15 + 12 + 56 + 27 + 30 = 143 = 13 \cdot 11 \equiv 0 \pmod{11}.
\end{align*}
Also ist $c$ ein Codewort. \checkmark

\medskip
{\small\textit{Hinweis:} Andere Formen von $H$ (z.\,B. systematisch $H = (A \mid E)$) beschreiben denselben Code.
Die Syndromwerte sehen dann anders aus, das Ergebnis (Fehlerstelle, Codewort) ist dasselbe.}
